% SHARD-ID: RC-03J-CCM-TENSOR-EXAMPLES
% SHARD-TITLE: CCM in Tensor-Network Terms IV: an Exact Elliptic-Curve Example
% SHARD-SUMMARY: Works through the retained F5 Pauli companion, its positive metric, and a three-length Weil matrix with an exact annihilating polynomial.
% SHARD-SUMMARY: Separates the spectral support recovered from counts from the physical Pauli letters and shows why the native Hilbert-Schmidt contraction gives the wrong Gram matrix.
% SHARD-KEYWORDS: CCM, Pauli channel, elliptic curve, F5, Gram matrix, metric, minimal polynomial
% SHARD-NOTES: none
% SHARD-SCRIPTS: scripts/ccm_tensor_network.py

\section{CCM in Tensor-Network Terms IV: an Exact Example}
\label{sec:ccm-tn-examples}

The net retained values of the graded nonbacktracking edge lift of
the six-letter Pauli example in \cref{sec:graded-ramanujan} are
$-1\pm2i$, of modulus $\sqrt5$; they are not eigenvalues of its
raw doubled adjacency transfer, whose spectrum is $6,-2,-2,-2$.
Here is the entire reconstruction on their two-dimensional companion
representative. It is small enough to exhibit every contraction exactly.

\begin{proposition}[The F5 retained companion and its Weil kernel]
\label{prop:ccm-tn-pauli}
\claimstatus{prop:ccm-tn-pauli}{proved}
For
\[
 E=\frac1{\sqrt5}\begin{pmatrix}0&-5\\1&-2\end{pmatrix},\qquad
 \TNmetric=\begin{pmatrix}1&-1\\-1&5\end{pmatrix},
\]
the metric is positive and $E^*\TNmetric E=\TNmetric$. The three-length
moment matrix is
\[
 W^{(3)}=\begin{pmatrix}
 2&-2/\sqrt5&-6/5\\
 -2/\sqrt5&2&-2/\sqrt5\\
 -6/5&-2/\sqrt5&2
 \end{pmatrix}.
\]
Its eigenvalues are $0,14/5,16/5$, and its kernel is spanned by
$\TNnull=(1,2/\sqrt5,1)^T$. Thus the kernel gives precisely
$E^2+(2/\sqrt5)E+I=0$. This matrix is the $G$-Hilbert--Schmidt Gram of
$I,E,E^2$, whereas it is not their ordinary Hilbert--Schmidt Gram.
\end{proposition}
\begin{proof}
\begin{enumerate}
\item[$\langle1\rangle$1.] The leading minors of $\TNmetric$ are $1,4$;
direct multiplication gives $E^*\TNmetric E=\TNmetric$.
\item[$\langle1\rangle$2.] Direct multiplication also gives the quadratic
relation, and $E$ is not scalar, so the polynomial is minimal. Taking
traces gives $t_0=2,t_1=-2/\sqrt5,t_2=-6/5$, hence the displayed matrix
by \cref{def:ccm-tn-moments}.
\item[$\langle1\rangle$3.] Its characteristic polynomial is
$s(s-14/5)(s-16/5)$ and $W^{(3)}\TNnull=0$, proving the rank and kernel.
\item[$\langle1\rangle$4.] Step 1 implies $E^\sharp=E^{-1}$, hence
$\Tr((E^j)^\sharp E^k)=t_{k-j}$. But
$\Tr(E^*E)=6\ne2=W^{(3)}_{11}$, disproving the native-metric identity.
\end{enumerate}
\end{proof}

\subsection{Reading the diagram from counts to the bond relation}

The unnormalised data are
\[
 N_k=1+5^k-\bigl((-1+2i)^k+(-1-2i)^k\bigr).
\]
Thus $N_0=0,N_1=8,N_2=32$. Remove the known even contribution with
the correct sign and scale the odd part:
\[
 t_k=5^{-k/2}(1+5^k-N_k).
\]
The resulting $W^{(3)}$ is built from three ring lengths. Its nullvector
specifies three open transfer networks whose coherent sum vanishes:
\[
 \underbrace{I}_{\text{length }0}
 +\frac2{\sqrt5}\underbrace{E}_{\text{length }1}
 +\underbrace{E^2}_{\text{length }2}=0.
\]
The kernel polynomial has roots $(-1\pm2i)/\sqrt5$; multiplication
by $z$ modulo that polynomial is the recovered two-dimensional moment
shift. The six Pauli letters supply these data through the graded
nonbacktracking edge lift and cancellation of its even/odd common
factors. The moment problem does not output those six letters.

The metric is equally tangible. Closing the open networks with
$A^\sharp=\TNmetric^{-1}A^*\TNmetric$ gives the correct overlaps;
closing with $A^*$ gives a different network. Similarity of transfer
matrices preserves ring traces while changing this uncorrected closure.

\subsection{Independent evidence}

The deterministic checker \texttt{scripts/ccm\_tensor\_network.py}
also checks the permutation $(1\ 2)(3\ 4\ 5)$, whose five-dimensional
permutation representation has only four distinct spectral values.
Its kernel polynomial is $(z^2-1)(z^2+z+1)$ and the recovered cyclic
space has dimension four. The repeated eigenvalue $1$ survives as a
weight, not as two independent moment vectors.
Complex phases test the clock conjugation convention; a Jordan block
tests the limit of spectral reconstruction; and the exact CCM toy in
the preceding shard tests the boundary correction.
These computations are finite evidence for the formulas, separate from
their proofs and from any infinite spectral convergence claim.
